% !PRINTED 505
\weberpage{505}

From this one obtains $\varkappa^2$ and $\varkappa'^2$ from the formulae [§ 54, (3)]
\[
\varkappa^2 = \frac{f_2(\omega)^8}{f(\omega)^8}, \quad \varkappa'^2 = \frac{f_1(\omega)^8}{f(\omega)^8},
\]
hence
\[
\begin{aligned}
\omega &= \sqrt{-2}\colon & \varkappa^2 &= (\sqrt{2} - 1)^2, \\
\omega &= \sqrt{-3}\colon & \varkappa^2 &= \frac{2 - \sqrt{3}}{4} = \frac{(1 - \sqrt{3})^2}{8}, \\
\omega &= \sqrt{-7}\colon & \varkappa^2 &= \frac{8 - 3\sqrt{7}}{16} = \frac{(3 - \sqrt{7})^2}{32}.
\end{aligned}
\]

\begin{center}
§ 135.\qquad \gesperrt{The class invariant $\varkappa^2$ and $\varkappa$.}
\end{center}

It has been shown in § 126 that, if the discriminant $D$ is divisible
by $4$, and the quadratic equation whose root is $\omega$ has odd
outer coefficients, then $f(\omega)^{24}$ is a class invariant
for the discriminant $D$.

But by § 54 and § 34 we have
\weq{1}{\varkappa^2 \varkappa'^2 = \frac{16}{f(\omega)^{24}}, \quad \varkappa^2 = \frac{f_2(\omega)^8}{f(\omega)^8} = \frac{16}{f(\omega)^8\, f_1(2\omega)^8}.}
If now $\omega$ satisfies the equation
\weq{2}{A\,\omega^2 + B\,\omega + C = 0}
with the discriminant
\[
D = B^2 - 4\,A\,C,
\]
then $\omega' = 2\,\omega$ satisfies the equation
\[
A\,\omega'^2 + 2\,B\,\omega' + 4\,C = 0
\]
with the discriminant $4\,D$, and by § 126, (13) therefore
$f_1(2\omega)^{24}$ is a class invariant for this discriminant. Thus
$\varkappa^2 \varkappa'^2$ and $\varkappa^6$ are rationally expressible by class invariants of the
discriminant $D$ and $4\,D$.

The identical equation
\weq{3}{\frac{\varkappa^2 \varkappa'^2 + \varkappa^4 \varkappa'^4 + 2\,\varkappa^6}{1 + \varkappa^6} = \varkappa^2}
shows that the same holds also for the modulus $\varkappa^2$.

Since $\varkappa^2$ passes under every linear substitution into a rational
(linear fractional) function of $\varkappa^2$, the assumption that the
coefficients $A$, $C$ are to be odd can now also be dropped, and we
have the theorem:

2. \gesperrt{If the discriminant $D \equiv 0 \pmod{4}$, then $\varkappa^2(\omega)$ is a class invariant of the discriminant $4\,D$.}
